% SPDX-License-Identifier: CC-BY-4.0
% Copyright (C) 2026 Fabian Januszewski, Abhijit and the Langlandscape team
% Stage 03, WP S03-M5: the GL_2 side (Betti only): amalgam, Mayer--Vietoris, Psi^*Psi in the normalisation of
% plan/implementation/stage-03.md sec. 0.2.8, the Betti level-raising module, Ihara for GL_2 (expected). Owner of if:s03-amalgam.
% Labels are fixed by plan/staged/stage-03-ribet-ihara.tex; new labels: lem:s03-groups, lem:s03-dcoset, lem:s03-Tclassical,
% lem:s03-PsiHecke, lem:s03-H1G.
\section{The \texorpdfstring{$\GL_2$}{GL2} side: the amalgam (Betti only)}\label{ss:s03-amalgam}
Let $(M,q)=1$ and $\Gamma:=\Gamma_0(M)$. Put $G:=\Gamma_0(M)[\frac1q]=\{\gamma\in\SL_2(\Z[\frac1q]):\gamma\equiv(\begin{smallmatrix}*&*\\0&*\end{smallmatrix})\bmod M\}$,
$\alpha'=\mathrm{diag}(q,1)$ and $\Gamma':=\alpha'^{-1}\Gamma\alpha'$. Then
$\alpha'^{-1}(\begin{smallmatrix}a&b\\c&d\end{smallmatrix})\alpha'=(\begin{smallmatrix}a&b/q\\qc&d\end{smallmatrix})$, so
$\Gamma\cap\Gamma'=\Gamma_0(Mq)$. Moreover $\alpha'^{-1}v_0=\alpha v_0=v_1$ in $\cT_q$, so the conventions agree with stage 02.
(Indeed $\alpha'^{-1}=q^{-1}\mathrm{diag}(1,q)$ and homotheties act trivially on $\cT_q$; $v_1=\alpha v_0$ by Def.~\ref{def:s02-tree}.)
We write $g^\iota:=\det(g)g^{-1}$ for the adjugate, so $\alpha'^\iota=\mathrm{diag}(1,q)=q\,\alpha'^{-1}$.

\begin{lemma}[elementary facts]\label{lem:s03-groups}\leavevmode
\begin{enumerate}[label=(\alph*)]
\item $\alpha'G\alpha'^{-1}=G$, $G\cap\SL_2(\Z)=\Gamma$, and $\Gamma\cap\alpha'\Gamma\alpha'^{-1}=\alpha'\Gamma_0(Mq)\alpha'^{-1}$.
\item $[\Gamma:\Gamma_0(Mq)]=[\Gamma':\Gamma_0(Mq)]=q+1$.
\item $\Gamma\alpha'^\iota\Gamma=\Gamma\alpha'\Gamma$; explicitly, if $a,b\in\Z$ with $aq-bM=1$, then $\gamma_1=\left(\begin{smallmatrix}a&b\\M&q\end{smallmatrix}\right)$ and
  $\gamma_2=\left(\begin{smallmatrix}qa&b\\M&1\end{smallmatrix}\right)$ lie in $\Gamma$ and $\alpha'^\iota=\gamma_1\alpha'\gamma_2^{-1}$.
\item The closure of $G$ in $\SL_2(\Qq)$ is $\SL_2(\Qq)$.
\end{enumerate}
\stproved.
\end{lemma}
\begin{proof}
(a) $\alpha'(\begin{smallmatrix}a&b\\c&d\end{smallmatrix})\alpha'^{-1}=(\begin{smallmatrix}a&qb\\c/q&d\end{smallmatrix})$ has entries in $\Z[\frac1q]$, and $c/q\equiv0\bmod M$ iff
$c\equiv0\bmod M$ ($q$ is a unit mod $M$); with the inverse computation this gives $\alpha'G\alpha'^{-1}=G$. An element of $G$ with integral entries is in
$\SL_2(\Z)$ with $M\mid c$, i.e.\ in $\Gamma$. The last identity is the conjugate by $\alpha'$ of $\Gamma'\cap\Gamma=\Gamma_0(Mq)$.

(b) Reduction mod $q$ maps $\Gamma$ onto $\SL_2(\F_q)$: $\Gamma$ contains $\left(\begin{smallmatrix}1&1\\0&1\end{smallmatrix}\right)$ and
$\left(\begin{smallmatrix}1&0\\M&1\end{smallmatrix}\right)$, whose reductions generate the upper unipotent group $N^+(\F_q)$ and, as $M\in\F_q^\times$, the lower
unipotent group $N^-(\F_q)$; these generate $\SL_2(\F_q)$ ($\SL_2$ of a field is generated by elementary matrices). $\Gamma_0(Mq)$ is the preimage of the upper triangular subgroup $B(\F_q)$, of index
$|\mathbb P^1(\F_q)|=q+1$ in $\SL_2(\F_q)$. Conjugating by $\alpha'$, $[\Gamma':\Gamma'\cap\Gamma]=[\Gamma:\Gamma\cap\alpha'\Gamma\alpha'^{-1}]$, and
$\Gamma\cap\alpha'\Gamma\alpha'^{-1}=\{\gamma\in\Gamma:q\mid b\}$ is the preimage of the lower triangular group, also of index $q+1$.

(c) $\det\gamma_1=aq-bM=1=\det\gamma_2$, and both have lower left entry $M$, so they lie in $\Gamma$. Then
$\gamma_1\alpha'=\left(\begin{smallmatrix}qa&b\\qM&q\end{smallmatrix}\right)=\mathrm{diag}(1,q)\gamma_2$.

(d) The closure $\bar G$ is a closed subgroup of $\SL_2(\Qq)$. $G$ contains $\left(\begin{smallmatrix}1&x\\0&1\end{smallmatrix}\right)$ and
$\left(\begin{smallmatrix}1&0\\Mx&1\end{smallmatrix}\right)$ for all $x\in\Z[\frac1q]$, which is dense in $\Qq$, as is $M\Z[\frac1q]$ ($M\in\Zq^\times$). So $\bar G$ contains both
unipotent subgroups $N^\pm(\Qq)$, which generate $\SL_2(\Qq)$ ($\SL_2$ of a field is generated by elementary matrices). This is hypothesis (D) below. (Alternatively it follows
from strong approximation for $\SL_2$ over $\Q$, \lit{voight\_2021\_quaternion-algebras}{Thm 28.2.3, p.~480 (proof p.~481)}, as in the stage tex.)
\end{proof}

\begin{proposition}[amalgam]\label{prop:s03-amalgam}
$G=\Gamma*_{\Gamma_0(Mq)}\Gamma'$. \stpmc: [Se, II.1.4, Thm~3, p.~79] \lit{serre\_1980\_trees}{II.1.4 Thm 3, p.~79} under
hypothesis (D) of [Se, II.1.4 Thm~2, p.~78] \lit{serre\_1980\_trees}{II.1.4 Thm 2, p.~78}: the closure of $G$ in $\SL_2(\Qq)$ is $\SL_2(\Qq)$
(Lemma~\ref{lem:s03-groups}(d), proved elementarily; the stage tex derives it from strong approximation for $\SL_2$ over $\Q$ [Vo, Thm~28.2.3], which is stated on
p.~480 (proof p.~481), not p.~479). For $M=1$ it is [Se, II.1.4, Cor.~2, p.~80; p.~81] \lit{serre\_1980\_trees}{II.1.4 Cor.~2, p.~80; p.~81} \stcited, whose second factor is
exactly $\{(\begin{smallmatrix}a&q^{-1}b\\qc&d\end{smallmatrix})\}=\Gamma'$. The vertex stabilisers are $\Gamma$ (of $v_0$) and $\Gamma'$ (of $v_1$).
\end{proposition}
\begin{proof}
$G\subseteq\SL_2(\Qq)$ acts on $\cT_q$ without inversion (Lemma~\ref{lem:s02-tree}(b)). Serre's Thm~3 (for $V=\Qq^2$ and the base edge $\{v_0,v_1\}$) says: under (D),
$G$ is the amalgam of the stabilisers of $v_0$ and $v_1$ along the stabiliser of the edge. These are: $G\cap\SL_2(\Zq)=G\cap\SL_2(\Z)=\Gamma$ (an element of
$\SL_2(\Z[\frac1q])$ integral at $q$ is integral; Lemma~\ref{lem:s03-groups}(a)); $\Stab(v_1)=\alpha'^{-1}\Stab(v_0)\alpha'$ in $\SL_2(\Qq)$, so
$G\cap\Stab(v_1)=\alpha'^{-1}(\alpha'G\alpha'^{-1}\cap\SL_2(\Zq))\alpha'=\alpha'^{-1}\Gamma\alpha'=\Gamma'$ by Lemma~\ref{lem:s03-groups}(a); and the edge stabiliser is
$\Gamma\cap\Gamma'=\Gamma_0(Mq)$ (no inversions). (D) is Lemma~\ref{lem:s03-groups}(d).
\end{proof}

Let $l$ be a prime with $l>k$, $l\ne q$, so $l\ge5$. Let $V=\Sym^{k-2}(\Z_{(l)}^2)$, $k\ge4$ even, be the symmetric power of the
standard (column) representation of $\GL_2(\Z_{(l)})\ni\alpha'$. Let $c:H^1(\Gamma,V)\xrightarrow\sim H^1(\Gamma',V)$,
$(cz)(\gamma')=\alpha'^{-1}z(\alpha'\gamma'\alpha'^{-1})$, be the conjugation isomorphism.
No condition $l\nmid M$ is needed (proof of Prop.~\ref{prop:s03-MV}). The hypotheses on the coefficients are therefore: $k\ge4$ even, $l$ prime,
$l>k$, $l\ne q$, $(M,q)=1$; they are listed again in \pathtt{if:s03-amalgam}.

\begin{proposition}[Mayer--Vietoris; seed \S6.4]\label{prop:s03-MV}
There is an exact sequence, with $\Psi(x,y)=\mathrm{res}\,x-\mathrm{res}\,(cy)$,
\[0\to H^1(G,V)\to H^1(\Gamma,V)^2\xrightarrow{\ \Psi\ }H^1(\Gamma_0(Mq),V)\to H^2(G,V)\to0.\]
\stpmc: the Mayer--Vietoris sequence of an amalgam with coefficients [Br, III.6, Exercise, p.~74]
\lit{brown\_1982\_cohomology-of-groups}{III.6 Exercise, p.~74} \stcited{} (the trivial-coefficient version is [Br, II.7 Cor.~7.7, p.~51]
\lit{brown\_1982\_cohomology-of-groups}{II Cor.~7.7, p.~51}), applied to Prop.~\ref{prop:s03-amalgam}. The end terms vanish for these reasons.
$H^0(\Gamma_0(Mq),V)=0$ for $k>2$: the image of $\Gamma_0(Mq)$ in $\SL_2(\F_l)$ is everything if $l\nmid M$, and $\Sym^{k-2}\F_l^2$ is irreducible and
non-trivial since $k-2<l$; if $l\mid M$ the image lies in the Borel subgroup, and the invariants vanish as well (proof). $H^2(\Gamma,V)=0$ because $\Gamma$ is virtually free with torsion of order prime to $l$. The same argument shows that $H^1(\Gamma,V)$ is torsion-free.
Moreover $H^1(G,V)=0$ (Lemma~\ref{lem:s03-H1G} below), so the sequence reads
\[0\to H^1(\Gamma,V)^2\xrightarrow{\ \Psi\ }H^1(\Gamma_0(Mq),V)\to H^2(G,V)\to0.\]
\end{proposition}
\begin{proof}
The Mayer--Vietoris sequence for $G=\Gamma*_{\Gamma_0(Mq)}\Gamma'$ reads
$\cdots\to H^0(\Gamma_0(Mq),V)\to H^1(G,V)\to H^1(\Gamma,V)\oplus H^1(\Gamma',V)\to H^1(\Gamma_0(Mq),V)\to H^2(G,V)\to H^2(\Gamma,V)\oplus H^2(\Gamma',V)$, the middle map
being a difference of restrictions (the sign convention does not affect exactness). Identifying $H^1(\Gamma',V)$ with $H^1(\Gamma,V)$ via $c$ gives $\Psi$.

\emph{$H^0(\Gamma_0(Mq),V)=0$.} Put $n=k-2$, so $1\le n\le l-3$. \emph{Case $l\nmid M$.} As in Lemma~\ref{lem:s03-groups}(b), $\Gamma_0(Mq)\to\SL_2(\F_l)$ is onto
($l\nmid Mq$). $\Sym^n\F_l^2$ is irreducible and non-trivial (proof of Remark~\ref{rem:s03-T6}(a), with $l$ for $p$; the action \eqref{eq:s01-Vk} and the column action
differ by transposition, which preserves irreducibility), so it has no non-zero invariants. \emph{Case $l\mid M$.} Then the image of $\Gamma_0(Mq)$ in $\SL_2(\F_l)$ lies in the Borel
subgroup. On $\Sym^n\F_l^2=\Sym^n(\F_le_1\oplus\F_le_2)$ the invariants of $\left(\begin{smallmatrix}1&1\\0&1\end{smallmatrix}\right)\in\Gamma_0(Mq)$ ($e_2\mapsto e_1+e_2$)
are the line $\F_le_1^n$ (as in the proof of Remark~\ref{rem:s03-T6}(a), since $n<l$). Choose $a\in\Z$ with $a\bmod l$ a generator of $\F_l^\times$ and $a\equiv1$ modulo
every prime $\ne l$ dividing $Mq$; then $(a,Mq)=1$, so $ad-bMq=1$ for some $b,d$, and $\gamma=\left(\begin{smallmatrix}a&b\\Mq&d\end{smallmatrix}\right)\in\Gamma_0(Mq)$ maps
$e_1\mapsto ae_1+Mqe_2\equiv ae_1$, hence acts on $e_1^n$ by $a^n\not\equiv1$ ($0<n<l-1$). So again there are no non-zero invariants.
(Since $\Gamma,\Gamma'\supseteq\Gamma_0(Mq)$, the same holds for them.) If $0\ne v\in V^{\Gamma_0(Mq)}$, divide by the largest power of $l$ to get $v\notin lV$; its reduction is a non-zero invariant of
$\Sym^n\F_l^2$, a contradiction.

\emph{$H^i(\Gamma,V)=0$ for $i\ge2$.} $\Gamma\subseteq\SL_2(\Z)$ acts on the tree of $\SL_2(\Z)$ with vertex stabilisers in conjugates of $\Z/6$, $\Z/4$ and edge stabilisers in
$\Z/2$ \lit{serre\_1980\_trees}{I.4.2(c), p.~35}, without inversion. These finite groups have order prime to $l$, so their higher cohomology with coefficients in the
$\Z_{(l)}$-module $V$ vanishes. The long exact sequence of the action (as in Prop.~\ref{prop:s02-coh}(a)) puts $H^i(\Gamma,V)$ ($i\ge2$) between
$\prod_eH^{i-1}(\Gamma_e,V)=0$ and $\prod_vH^i(\Gamma_v,V)=0$ (finitely many orbits, $\Gamma$ of finite index). The same holds for $\Gamma'$ (conjugate) and $\Gamma_0(Mq)$.

\emph{Torsion-freeness.} From $0\to V\xrightarrow{l}V\to V/l\to0$: $H^1(\Gamma,V)[l]$ is the cokernel of $H^0(\Gamma,V)\to H^0(\Gamma,V/l)$, and
$H^0(\Gamma,V/l)=(\Sym^n\F_l^2)^\Gamma=0$ as above. As $H^1(\Gamma,V)$ is a $\Z_{(l)}$-module, it is torsion-free. (Also $H^1(\Gamma,V)/l\cong H^1(\Gamma,V/l)$, as
$H^2(\Gamma,V)=0$.)
\end{proof}

\subsection*{Double-coset operators and \texorpdfstring{$\Psi^*\Psi$}{Psi*Psi}}
For a subgroup $H\subseteq\GL_2(\Z_{(l)})$ and $g\in\GL_2(\Z_{(l)})$ let $c_g:H^*(H,V)\to H^*(gHg^{-1},V)$ be the conjugation map, induced by
$(h,v)\mapsto(ghg^{-1},gv)$; on $1$-cocycles $(c_gz)(h')=g\,z(g^{-1}h'g)$
\lit{brown\_1982\_cohomology-of-groups}{III.8, p.~80}. Then $c=c_{\alpha'^{-1}}$ and $c^{-1}=c_{\alpha'}$. For $g\in\GL_2(\Z_{(l)})$ commensurating $\Gamma$ (i.e.\ $\Gamma\cap g\Gamma g^{-1}$ of finite index in $\Gamma$ and in $g\Gamma g^{-1}$) put
\[T_g:=\mathrm{cor}^{\Gamma}_{\Gamma\cap g\Gamma g^{-1}}\circ c_g\circ\mathrm{res}^{\Gamma}_{\Gamma\cap g^{-1}\Gamma g}:\ H^*(\Gamma,V)\to H^*(\Gamma,V),\]
the double-coset operator of $\Gamma g\Gamma$ acting on $V$ through $g$.

\begin{lemma}[double cosets]\label{lem:s03-dcoset}\leavevmode
\begin{enumerate}[label=(\alph*)]
\item $T_{\gamma_1g\gamma_2}=T_g$ for $\gamma_1,\gamma_2\in\Gamma$; so $T_g$ depends only on $\Gamma g\Gamma$ (and on the action of its elements on $V$).
\item $T_{\lambda g}=\lambda^{k-2}T_g$ for a scalar $\lambda\in\Z_{(l)}^\times$.
\item $\mathrm{cor}^\Gamma_H\circ\mathrm{res}^\Gamma_H=[\Gamma:H]$.
\end{enumerate}
\stpmc{} (formal properties of $\mathrm{res},\mathrm{cor},c_g$ cited from Brown).
\end{lemma}
\begin{proof}
(a) $c_{g}$ is functorial, $c_{gh}=c_gc_h$, and commutes with restriction and corestriction ($c_g\circ\mathrm{res}^H_K=\mathrm{res}^{gHg^{-1}}_{gKg^{-1}}\circ c_g$,
likewise for $\mathrm{cor}$) \lit{brown\_1982\_cohomology-of-groups}{III.9, Exercise 1, p.~83}; and $c_\gamma=\mathrm{id}$ on $H^*(\Gamma,V)$ for $\gamma\in\Gamma$
\lit{brown\_1982\_cohomology-of-groups}{III.8 Prop.~(8.3), p.~80}. With $H:=\Gamma\cap g\Gamma g^{-1}$, $H':=\Gamma\cap g^{-1}\Gamma g$ and $g'=\gamma_1g\gamma_2$:
$\Gamma\cap g'\Gamma g'^{-1}=\gamma_1H\gamma_1^{-1}$ and $\Gamma\cap g'^{-1}\Gamma g'=\gamma_2^{-1}H'\gamma_2$, so
$T_{g'}=\mathrm{cor}^\Gamma_{\gamma_1H\gamma_1^{-1}}c_{\gamma_1}c_gc_{\gamma_2}\mathrm{res}^\Gamma_{\gamma_2^{-1}H'\gamma_2}=c_{\gamma_1}\,\mathrm{cor}^\Gamma_Hc_g\mathrm{res}^\Gamma_{H'}\,c_{\gamma_2}=T_g$.
(b) The groups are unchanged, and $c_{\lambda g}=\rho(\lambda)c_g$ with $\rho(\lambda)=\lambda^{k-2}$ on $\Sym^{k-2}$. (c) \lit{brown\_1982\_cohomology-of-groups}{III.9
Prop.~(9.5)(ii), p.~82}.
\end{proof}

\begin{lemma}[$T_{\alpha'}$ is the classical $T_q$]\label{lem:s03-Tclassical}
Let $n=k-2$ and let $\iota$ be the adjugate, $g^\iota=\det(g)g^{-1}$; it is an anti-automorphism of $M_2$ with $\gamma^\iota=\gamma^{-1}$ on $\SL_2$, so $\Gamma^\iota=\Gamma$.
\begin{enumerate}[label=(\alph*)]
\item \emph{Cocycle formula.} Let $g\in\GL_2(\Z_{(l)})$ commensurate $\Gamma$, $H:=\Gamma\cap g\Gamma g^{-1}$, $\Gamma=\coprod_ir_iH$, and $\beta_i:=r_ig$. Then
  $\Gamma g\Gamma=\coprod_i\beta_i\Gamma$, and on $1$-cocycles
  \[(T_gz)(\gamma)=\sum_i\beta_i\,z\big(\beta_i^{-1}\gamma\beta_{j(i)}\big),\qquad\text{where }\gamma\beta_{j(i)}\in\beta_i\Gamma.\]
\item \emph{Hida's coefficients.} Let $L(n)$ be the module of homogeneous polynomials of degree $n$ in $X,Y$ over $\Z_{(l)}$, with Hida's action
  $(\gamma P)(x)=P(\gamma^\iota x)$, $x={}^{t}(X,Y)$, of the monoid $M_2(\Z_{(l)})\cap\GL_2(\Q)$
  \lit{hida\_1993\_elementary-theory-of-l-functions-and-eisenstein-series}{\S6.1, p.~165}. For a column vector $v$ put $\ell_v(x):=\det(v,x)={}^{t}v\,w\,x$,
  $w=\left(\begin{smallmatrix}0&1\\-1&0\end{smallmatrix}\right)$. Then $v\mapsto\ell_v$ induces an isomorphism $\Sym^n(\Z_{(l)}^2)\xrightarrow\sim L(n)$ of modules over
  this monoid (standard column action on the left); a scalar $t$ acts by $t^n$ on both.
\item \emph{Hida's operator.} For a decomposition $\Gamma\alpha'\Gamma=\coprod_i\Gamma\alpha_i$ and $\alpha_i\gamma=\gamma_i\alpha_{j}$ ($\gamma,\gamma_i\in\Gamma$),
  Hida's operator $[\Gamma\alpha'\Gamma]$ on $1$-cocycles is $v(\gamma)=\sum_i\alpha_i^\iota u(\gamma_i)$
  \lit{hida\_1993\_elementary-theory-of-l-functions-and-eisenstein-series}{\S6.3, pp.~175--176}. Under (b) it equals $T_{\alpha'}$ on $H^1(\Gamma,V)$.
\item \emph{Classical $T_q$.} Hence $T_q:=T_{\alpha'}$ is the classical Hecke operator $f\mapsto f|[\Gamma\alpha'\Gamma]=\sum_if|\alpha_i$, with
  $f|\alpha=\det(\alpha)^{k-1}f(\alpha z)j(\alpha,z)^{-k}$ (eigenvalue $a_q$ on a normalised eigenform), under the Eichler--Shimura isomorphism
  $S_k(\Gamma)\oplus\overline{S_k(\Gamma)}\cong H^1_P(\Gamma,V\otimes\C)$, which is Hecke equivariant
  \lit{hida\_1993\_elementary-theory-of-l-functions-and-eisenstein-series}{\S6.3, pp.~176--177, Thm.~1 p.~177}.
\item \emph{Eisenstein part.} Let $k\ge4$ even, $V_\C:=V\otimes\C\cong L(n;\C)$ by (b). For $f\in M_k(\Gamma)$ and $z\in\mathfrak H$ put
  $\varphi_z(f)(\gamma):=\int_z^{\gamma z}f(w)(X-wY)^n\,dw$. Then $\varphi_z(f)$ is a $1$-cocycle whose class $\varphi(f)\in H^1(\Gamma,V_\C)$ does not depend on $z$,
  $\varphi(f|[\Gamma\alpha'\Gamma])=T_q\varphi(f)$ for all $f\in M_k(\Gamma)$, and
  \[H^1(\Gamma,V_\C)=H^1_P(\Gamma,V_\C)\oplus\varphi(E_k(\Gamma)),\]
  with both summands $T_q$-stable and $\varphi:E_k(\Gamma)\to\varphi(E_k(\Gamma))$ a $T_q$-equivariant isomorphism ($E_k(\Gamma)$ the Eisenstein subspace
  \lit{stein\_2007\_modular-forms-a-computational-approach\_draft}{\S5.1, p.~83}). Hence the generalised eigenvalues of $T_q$ on $H^1(\Gamma,V_\C)$ are those of the
  classical $T_q$ on $S_k(\Gamma)$, on $\overline{S_k(\Gamma)}$ and on $E_k(\Gamma)$.
\end{enumerate}
\emph{Status.} (a)--(c) \stproved{} (from the cochain description of $\mathrm{cor}$ in Brown); (d) \stpmc{} (Hida's Thm.~1, \stcited); (e) \stpmc{}
(proof below; it uses Hida's computation \S6.3, pp.~176--177, for all of $M_k$, Hida's Thm.~1 for $H^1_P$, and only the dimension count of Stein's (A.2.4)).
Stein states without proof that $T_{\mathrm{diag}(1,p)}$ on all of $H^1(\SL_2(\Z);M_{k-2})$ acts as the classical $T_p$
\lit{stein\_2007\_modular-forms-a-computational-approach\_draft}{Example A.17, p.~226} \stcited. \stcomp{}, \pathtt{code/results/s03/C3\_rerun.json}
(S03-H1, HACS exact $\mathrm{char}_\Z(T_q\mid H^1(\Gamma_0(M);\Sym^{k-2}\Z^2))$; provenance: commit \texttt{05f7fe72}, script \pathtt{code/hacs/jl/s03\_h1\_rerun.py},
hacs 0.2.1): for all $14$ cases of the file ($M=1$, $k=12$; $M=11$, $k\in\{2,4,12\}$; $M=5$, $k=4$; $q\in\{2,3,5\}$, $q=5$ omitted for $M=5$), this polynomial (field \pathtt{hacs\_charpoly\_H1\_high\_to\_low}) factors, by an exact check of the JSON fields, as
$\mathrm{char}(T_q\mid S_k(\Gamma_0(M)))^2\cdot(X-1-q^{k-1})^{e}$ with $e=\dim E_k(\Gamma_0(M))$ ($e=1$ for $M=1$ and for $(M,k)=(11,2)$, $e=2$ otherwise), e.g.\ $(X+24)^2(X-2049)$ for
$(M,k,q)=(1,12,2)$, $(X-252)^2(X-177148)$ for $(1,12,3)$, $(X-4830)^2(X-48828126)$ for $(1,12,5)$, $(X+4)^2(X-9)^2$ for $(5,4,2)$.
\end{lemma}
\begin{proof}
(a) The stabiliser of $g\Gamma$ in $\Gamma$ is $H$, so $\Gamma g\Gamma=\coprod_ir_ig\Gamma$. For $\mathrm{cor}^\Gamma_H$ on cochains, Brown extends an $H$-cochain along the map
$\rho:\Gamma\to H$ of left $H$-sets sending each chosen representative of $H\backslash\Gamma$ to $1$, and then applies $f\mapsto\sum_{r\in\Gamma/H}rf$
\lit{brown\_1982\_cohomology-of-groups}{III.9 (C), (D), pp.~81--82}. With the representatives $r_i^{-1}$ of $H\backslash\Gamma$ and the inhomogeneous cocycle $y$ of the
homogeneous cochain $(h_0,h_1)\mapsto h_0y(h_0^{-1}h_1)$, this gives $(\mathrm{cor}\,y)(\gamma)=\sum_ir_i\,y(r_i^{-1}\gamma r_{j(i)})$, where $\gamma r_{j(i)}\in r_iH$
(indeed $\rho(r_i^{-1})=1$ and $\rho(r_i^{-1}\gamma)=r_i^{-1}\gamma r_{j}$ for the $j$ with $Hr_i^{-1}\gamma=Hr_j^{-1}$). Apply this to $y=c_g\,\mathrm{res}\,z$,
$y(h)=g\,z(g^{-1}hg)$: the summand is $r_ig\,z(g^{-1}r_i^{-1}\gamma r_jg)=\beta_iz(\beta_i^{-1}\gamma\beta_j)$, and $\gamma r_j\in r_iH$ iff $\gamma\beta_j\in\beta_ig^{-1}Hg\subseteq\beta_i\Gamma$;
as the $\beta_i\Gamma$ are distinct, this determines $j(i)$.

(b) For any $2\times2$ matrix $\gamma$, ${}^t\gamma\,w=w\,\gamma^\iota$, so $\ell_v(\gamma^\iota x)={}^tv\,{}^t\gamma\,w\,x=\ell_{\gamma v}(x)$, i.e.\ $\gamma\ell_v=\ell_{\gamma v}$. So
$v\mapsto\ell_v$ is an isomorphism of $\Z_{(l)}^2$ onto $L(1)$ ($w\in\SL_2(\Z)$), compatible with the actions, and $L(n)=\Sym^nL(1)$.

(c) Hida's operator does not depend on the decomposition: replacing $\alpha_i$ by $\delta_i\alpha_i$ ($\delta_i\in\Gamma$) replaces $\gamma_i$ by $\delta_i\gamma_i\delta_j^{-1}$, and with
$\alpha_i^\iota\gamma_i=\gamma\alpha_j^\iota$ (the image of $\alpha_i\gamma=\gamma_i\alpha_j$ under $\iota$) and the cocycle rule one finds $v'=v-(\gamma-1)x$,
$x=\sum_i\alpha_i^\iota\delta_i^{-1}u(\delta_i)$. Now take $g:=\alpha'^\iota$ in (a) and $\alpha_i:=\beta_i^\iota$. Applying $\iota$ to $\Gamma g\Gamma=\coprod_i\beta_i\Gamma$ gives
$\Gamma\alpha'\Gamma=\coprod_i\Gamma\alpha_i$, and applying it to $\alpha_i\gamma=\gamma_i\alpha_j$ gives $\gamma^{-1}\beta_i=\beta_j\gamma_i^{-1}$, i.e.\ $\gamma\beta_j=\beta_i\gamma_i$, so
$j=j(i)$ and $\gamma_i=\beta_i^{-1}\gamma\beta_{j(i)}$. Hence $v(\gamma)=\sum_i\beta_iu(\beta_i^{-1}\gamma\beta_{j(i)})=(T_{\alpha'^\iota}u)(\gamma)$ by (a), where $\beta_i$ acts on
$L(n)$ by Hida's action, which is the standard action under (b). Finally $T_{\alpha'^\iota}=T_{\alpha'}$ by Lemma~\ref{lem:s03-dcoset}(a) and Lemma~\ref{lem:s03-groups}(c).

(d) Hida shows $\varphi(f|[\Gamma\alpha\Gamma])=\varphi(f)|[\Gamma\alpha\Gamma]$ with exactly the operator of (c) on cocycles and $f|[\Gamma\alpha\Gamma]=\sum_if|\alpha_i$
\lit{hida\_1993\_elementary-theory-of-l-functions-and-eisenstein-series}{\S6.3, pp.~176--177}, and states the result for $\Gamma_0(N)$ as Thm.~1, p.~177 (trivial
character here). The classical $T_q$ ($q\nmid M$) is $f\mapsto f|[\Gamma\,\mathrm{diag}(1,q)\,\Gamma]$ in this normalisation, and $\Gamma\,\mathrm{diag}(1,q)\,\Gamma=\Gamma\alpha'^\iota\Gamma=\Gamma\alpha'\Gamma$
(Lemma~\ref{lem:s03-groups}(c)). \stknown

(e) \emph{Cocycle.} Put $\omega(f):=f(w)(X-wY)^n\,dw$, holomorphic on $\mathfrak H$. For $\alpha\in\GL_2^+(\Q)$, $\alpha^\iota\,(\omega(f)\circ\alpha)=\omega(f|\alpha)$ with
$f|\alpha=\det(\alpha)^{k-1}f(\alpha w)j(\alpha,w)^{-k}$ \lit{hida\_1993\_elementary-theory-of-l-functions-and-eisenstein-series}{\S6.3, p.~176}; for $\gamma\in\Gamma$ this is
$\omega(f)\circ\gamma=\gamma\,\omega(f)$. As $\mathfrak H$ is simply connected, the integral does not depend on the path, and Hida's computation
\lit{hida\_1993\_elementary-theory-of-l-functions-and-eisenstein-series}{\S6.2, pp.~167--168} gives the cocycle rule and
$\varphi_z(f)-\varphi_{z'}(f)=\big(\gamma\mapsto(1-\gamma)\int_{z'}^{z}\omega(f)\big)$, a coboundary. Hida states this for cusp forms, but for base points $z,z'\in\mathfrak H$ only
holomorphy and the transformation law are used (the decay at the cusps is needed only for base points at cusps). This is the $\C$-linear map of the proof of Hida's
Thm.~1, p.~177, extended from $S_k$ to $M_k$.

\emph{Hecke compatibility.} Hida's computation \lit{hida\_1993\_elementary-theory-of-l-functions-and-eisenstein-series}{\S6.3, pp.~176--177} of
$\varphi_z(f|[\Gamma\alpha\Gamma])$ fixes $z\in\mathfrak H$ and uses only the identity $\alpha^\iota(\omega(f)\circ\alpha)=\omega(f|\alpha)$, the primitive
$F(w)=\int^w\omega(f)+a$ and the coset relations $\alpha_i\gamma=\gamma_i\alpha_j$; cuspidality of $f$ is not used. So $\varphi(f|[\Gamma\alpha'\Gamma])=\varphi(f)|[\Gamma\alpha'\Gamma]$
for every $f\in M_k(\Gamma)$, and the right-hand side is $T_q\varphi(f)$ by (c), (d) (over $\C$ by flat base change, Hida \S6.2 (1a), p.~168).

\emph{Boundary.} Let $s$ be a cusp, $\sigma\in\SL_2(\Z)$ with $\sigma\infty=s$, $h\ge1$ with $t_h:=\left(\begin{smallmatrix}1&h\\0&1\end{smallmatrix}\right)\in\sigma^{-1}\Gamma\sigma$, and
$\pi:=\sigma t_h\sigma^{-1}\in\Gamma$. Substituting $w=\sigma w'$ and using $\omega(f)\circ\sigma=\sigma\,\omega(f|\sigma)$:
$\varphi_z(f)(\pi)=\sigma\,\varphi_{z'}(f|\sigma)(t_h)$, $z'=\sigma^{-1}z$. On $L(n;\C)$, $(t_hP)(X,Y)=P(X-hY,Y)$ (Hida's action, $t_h^\iota=t_{-h}$), so
$(t_h-1)L(n;\C)$ lies in the span of the $X^iY^{n-i}$ with $i<n$; and $(\pi-1)L(n;\C)=\sigma(t_h-1)L(n;\C)$. The coefficient of $X^n$ in
$\varphi_{z'}(f|\sigma)(t_h)$ is $\int_{z'}^{z'+h}(f|\sigma)(w)\,dw=h\,a_0(f|\sigma)$, where $a_0$ is the constant term of the expansion of $f|\sigma$ in
$e^{2\pi iw/h}$. Now let $f\in E_k(\Gamma)$ with $\varphi(f)\in H^1_P(\Gamma,V_\C)$, i.e.\ $\varphi_z(f)(\pi)\in(\pi-1)L(n;\C)$ for every parabolic $\pi\in\Gamma$ (Hida's
definition, \S6.2, p.~168; the condition does not depend on the representing cocycle). Then $a_0(f|\sigma)=0$ for every cusp, so $f\in S_k(\Gamma)\cap E_k(\Gamma)=0$
(\lit{stein\_2007\_modular-forms-a-computational-approach\_draft}{\S5.1, p.~83}: $M_k=E_k\oplus S_k$, $S_k$ the forms vanishing at all cusps). So
$E_k(\Gamma)\to H^1(\Gamma,V_\C)/H^1_P(\Gamma,V_\C)$ is injective.

\emph{Dimensions.} By Hida's Thm.~1, $\dim H^1_P(\Gamma,V_\C)=2\dim S_k(\Gamma)$. By Stein's (A.2.4)
\lit{stein\_2007\_modular-forms-a-computational-approach\_draft}{App.~A.2.1, (A.2.4), p.~205}, of which only the dimension is used, $\dim H^1(\Gamma,V_\C)=2\dim S_k(\Gamma)+\dim E_k(\Gamma)$; there
the coefficient module $M_{k-2}$ with $\gamma\cdot p(x,y)=p(dx-by,-cx+ay)$ (A.2.2) is $L(n;\C)$ with Hida's action, and the cohomology of $Y_0(N)$ with complex coefficients is
group cohomology also for $\Gamma$ with torsion (Stein, p.~211). Hence $\varphi(E_k(\Gamma))$ is a complement of $H^1_P$.

\emph{Stability.} $T_q$ preserves $H^1_P$ (Hida \S6.3, p.~176) and $E_k(\Gamma)$ (the decomposition $M_k=E_k\oplus S_k$ is one of Hecke modules, Stein \S5.1, p.~83), so by the
Hecke compatibility $\varphi(E_k(\Gamma))$ is $T_q$-stable and $\varphi|_{E_k}$ is $T_q$-equivariant. On $H^1_P$ the eigenvalues are those on $S_k\oplus\overline{S_k}$ by (d).
\end{proof}

\begin{proposition}[$\Psi^*\Psi$]\label{prop:s03-PsiPsi}
Put $\Psi^*(z):=\big(\mathrm{cor}^\Gamma_{\Gamma_0(Mq)}z,\ -c^{-1}\mathrm{cor}^{\Gamma'}_{\Gamma_0(Mq)}z\big)$. Then
\[\Psi^*\Psi=\begin{pmatrix}q+1&-T'_q\\-T''_q&q+1\end{pmatrix}\ \text{ on }H^1(\Gamma,V)^2,\qquad
T'_q=\mathrm{cor}\circ\mathrm{res}\circ c,\quad T''_q=c^{-1}\circ\mathrm{cor}'\circ\mathrm{res}.\]
Here $T'_q=T_{\alpha'^{-1}}$ is the operator of the double coset $\Gamma\alpha'^{-1}\Gamma=q^{-1}\Gamma\alpha'^\iota\Gamma$ (acting through its elements), and
$T''_q=T_{\alpha'}$ is that of $\Gamma\alpha'\Gamma$. For $V=\Sym^{k-2}$ of the standard (column) representation and $T_q:=T_{\alpha'}$, the classical
operator (Lemma~\ref{lem:s03-Tclassical}),
\[T'_q=q^{2-k}T_q,\qquad T''_q=T_q,\qquad \Psi^*\Psi=\begin{pmatrix}q+1&-q^{2-k}T_q\\-T_q&q+1\end{pmatrix},\]
\[\det(\Psi^*\Psi)=(q+1)^2-q^{2-k}T_q^2=-q^{2-k}\big(T_q^2-q^{k-2}(q+1)^2\big),\]
which is Ribet's level-raising condition up to a unit. \stpmc{} (formal properties of $\mathrm{res},\mathrm{cor},c_g$ cited from Brown, Lemma~\ref{lem:s03-dcoset};
the identification of $T_q$ with the classical operator is Lemma~\ref{lem:s03-Tclassical}(d); the stage tex has \sttoprove). The determinant is \stcomp{}
(see after the table below).
\end{proposition}

\paragraph{The normalisation table (stage-03.md \S0.2.8).} With $V=\Sym^{k-2}$ of the standard column representation and
$T_q:=T_{\alpha'}$ (Lemma~\ref{lem:s03-dcoset}; the double coset $\Gamma\alpha'\Gamma=\Gamma\alpha'^\iota\Gamma$ acting through its \emph{integral} elements; the classical
$T_q$ by Lemma~\ref{lem:s03-Tclassical}):
\begin{center}\small
\begin{tabular}{lll}\toprule
entry of $\Psi^*\Psi$ & operator & value\\\midrule
$(1,1)$, $(2,2)$ & $\mathrm{cor}\circ\mathrm{res}$, $c^{-1}\mathrm{cor}'\mathrm{res}'c$ & $q+1$\\
$(1,2)$ & $-T'_q=-\mathrm{cor}\circ\mathrm{res}\circ c=-T_{\alpha'^{-1}}$ (through $\alpha'^{-1}$) & $-q^{2-k}T_q$\\
$(2,1)$ & $-T''_q=-c^{-1}\circ\mathrm{cor}'\circ\mathrm{res}=-T_{\alpha'}$ (through $\alpha'$) & $-T_q$\\\bottomrule
\end{tabular}
\end{center}
So $\Psi^*\Psi=\left(\begin{smallmatrix}q+1&-q^{2-k}T_q\\-T_q&q+1\end{smallmatrix}\right)$: \textbf{the $(1,2)$ entry carries $q^{2-k}$}. This is the table of
\S0.2.8 of \pathtt{plan/implementation/stage-03.md}; the proof below confirms it (no correction).
\stcomp{}, \pathtt{code/results/s03/betti\_det.json} (S03-C5, PARI 2.15.4 on $S_k(\Gamma_0(M))$; provenance: commit \texttt{0f3bd5e1}, script
\pathtt{code/python/langlandscape/ribet/betti.py}): the file records this table (slot of $q^{2-k}$: $(1,2)$) and, for its $121$ cases ($M\in\{1,2,3,5,7,11\}$,
$k\in\{2,4,\dots,12\}$, primes $q\le13$, $q\nmid M$), $\det=(q+1)^2-q^{2-k}\lambda^2$ on each $T_q$-eigenvalue $\lambda$ of $S_k(\Gamma_0(M))$; it never vanishes
(\pathtt{vanishing\_triples\_M\_k\_q} is empty). E.g.\ $(M,k,q)=(1,12,2)$: $\lambda=-24$, $N(q^{k-2}(q+1)^2-\lambda^2)=8640=2^6\cdot3^3\cdot5$; $(1,12,3)$: $\lambda=252$,
$881280=2^7\cdot3^4\cdot5\cdot17$. \pathtt{code/results/s03/C3\_rerun.json} (S03-H1, HACS, commit \texttt{05f7fe72}) agrees on all $14$ overlapping cases, e.g.\
$\det=135/16$ for $(1,12,2)$. Which off-diagonal slot carries $q^{2-k}$ is \emph{not} tested numerically (field \pathtt{off\_diagonal\_assignment}: ``untested by
PARI''); it rests on the proof, and $\det$ and $\coker$ do not depend on it.

\noindent\emph{Other choices (recorded, not used).} The seed's matrix $\left(\begin{smallmatrix}q+1&-T_q\\-q^{2-k}T_q&q+1\end{smallmatrix}\right)$ is obtained by
exchanging $\alpha'$ and $\alpha'^{-1}$, i.e.\ by using $\Gamma'':=\alpha'\Gamma\alpha'^{-1}$ and $c_{\alpha'}$ in place of $\Gamma'$ and $c$ (then $T'_q=T_{\alpha'}$,
$T''_q=T_{\alpha'^{-1}}$). Likewise, which off-diagonal entry carries $q^{2-k}$ depends on whether $V$ or $V^\vee$ is the coefficient system. All choices
have the same $\det$ and isomorphic cokernels (end of the proof below).

\begin{proof}
\emph{The entries.} $\Psi^*\Psi(x,y)=\big(\mathrm{cor}\,\mathrm{res}\,x-\mathrm{cor}\,\mathrm{res}'\,cy,\ -c^{-1}\mathrm{cor}'\mathrm{res}\,x+c^{-1}\mathrm{cor}'\mathrm{res}'cy\big)$, where
$\mathrm{res},\mathrm{cor}$ are between $\Gamma$ and $\Gamma_0(Mq)$, and $\mathrm{res}',\mathrm{cor}'$ between $\Gamma'$ and $\Gamma_0(Mq)$. By Lemma~\ref{lem:s03-dcoset}(c) and
Lemma~\ref{lem:s03-groups}(b), $\mathrm{cor}\,\mathrm{res}=q+1$ and $\mathrm{cor}'\mathrm{res}'=q+1$, so the diagonal entries are $q+1$. The off-diagonal entries are
$-T'_q$ and $-T''_q$ by definition.

\emph{$T'_q=T_{\alpha'^{-1}}$.} With $g=\alpha'^{-1}$: $\Gamma\cap g\Gamma g^{-1}=\Gamma\cap\Gamma'=\Gamma_0(Mq)$ and $\Gamma\cap g^{-1}\Gamma g=\alpha'\Gamma_0(Mq)\alpha'^{-1}$
(Lemma~\ref{lem:s03-groups}(a)). Since $c=c_{\alpha'^{-1}}$ commutes with restriction,
$\mathrm{res}^{\Gamma'}_{\Gamma_0(Mq)}\circ c=c_{\alpha'^{-1}}\circ\mathrm{res}^\Gamma_{\alpha'\Gamma_0(Mq)\alpha'^{-1}}$, hence $T'_q=\mathrm{cor}^\Gamma_{\Gamma_0(Mq)}c_{\alpha'^{-1}}\mathrm{res}^\Gamma_{\alpha'\Gamma_0(Mq)\alpha'^{-1}}=T_{\alpha'^{-1}}$.

\emph{$T''_q=T_{\alpha'}$.} With $g=\alpha'$: $\Gamma\cap g\Gamma g^{-1}=\alpha'\Gamma_0(Mq)\alpha'^{-1}$, $\Gamma\cap g^{-1}\Gamma g=\Gamma_0(Mq)$, and
$c_{\alpha'}\circ\mathrm{cor}^{\Gamma'}_{\Gamma_0(Mq)}=\mathrm{cor}^{\Gamma}_{\alpha'\Gamma_0(Mq)\alpha'^{-1}}\circ c_{\alpha'}$. So $T''_q=T_{\alpha'}=T_q$.

\emph{The power of $q$.} $\alpha'^{-1}=q^{-1}\alpha'^\iota$, so by Lemma~\ref{lem:s03-dcoset}(b),(a) and Lemma~\ref{lem:s03-groups}(c),
$T_{\alpha'^{-1}}=q^{2-k}T_{\alpha'^\iota}=q^{2-k}T_{\alpha'}$. Hence $T'_q=q^{2-k}T_q$ and $T''_q=T_q$.

\emph{Determinant.} The entries lie in the commutative algebra $\Z_{(l)}[T_q]$ ($q$ is a unit in $\Z_{(l)}$), so
$\det(\Psi^*\Psi)=(q+1)^2-T'_qT''_q=(q+1)^2-q^{2-k}T_q^2=-q^{2-k}\big(T_q^2-q^{k-2}(q+1)^2\big)$. The product $T'_qT''_q=q^{2-k}T_q^2$ is the same for the other choice
(the factors are exchanged), and replacing $V$ by $V^\vee\cong V\otimes{\det}^{2-k}$ (an isomorphism over $\Z_{(l)}$ because $l>k-2$) multiplies $T_{\alpha'}$ by $q^{2-k}$
and $T_{\alpha'^{-1}}$ by $q^{k-2}$, which again leaves the product unchanged.

\emph{Cokernel.} Put $D:=\mathrm{diag}(1,q^{2-k})$, acting on $H^1(\Gamma,V)^2$; it is an automorphism ($q\in\Z_{(l)}^\times$) commuting with every Hecke operator
acting diagonally. Then
\[D\begin{pmatrix}q+1&-q^{2-k}T_q\\-T_q&q+1\end{pmatrix}D^{-1}=\begin{pmatrix}q+1&-T_q\\-q^{2-k}T_q&q+1\end{pmatrix},\]
so $D$ induces an isomorphism of Hecke modules between the cokernels of our matrix and of the seed's matrix. For $V^\vee$ the matrix is again of one of these two
shapes (with $T_q$ replaced by the operator on $V^\vee$), by the rescaling just described. So $\det(\Psi^*\Psi)$ and $\coker(\Psi^*\Psi)$ up to isomorphism do not
depend on the normalisation (D4).
\end{proof}

\begin{lemma}[Hecke equivariance of $\Psi$]\label{lem:s03-PsiHecke}
Let $p\nmid Mql$ be prime and $\beta:=\mathrm{diag}(p,1)\in\GL_2(\Z_{(l)})$. For $\Delta\in\{\Gamma,\Gamma',\Gamma_0(Mq)\}$ let $T_p^\Delta$ be the operator
$T_\beta$ defined as above with $\Delta$ in place of $\Gamma$. Then $\mathrm{res}^\Gamma_{\Gamma_0(Mq)}T_p^\Gamma=T_p^{\Gamma_0(Mq)}\mathrm{res}^\Gamma_{\Gamma_0(Mq)}$, likewise for
$\Gamma'$, and $c\,T_p^\Gamma=T_p^{\Gamma'}c$. Hence $\Psi\circ(T^\Gamma_p\oplus T^\Gamma_p)=T^{\Gamma_0(Mq)}_p\circ\Psi$. Let $\T^{(Mql)}$ be the $\Z_{(l)}$-algebra generated by the
pairs $(T_p^\Gamma\oplus T_p^\Gamma,\,T_p^{\Gamma_0(Mq)})$, $p\nmid Mql$, acting on $H^1(\Gamma,V)^2\oplus H^1(\Gamma_0(Mq),V)$; then $\Psi$ is $\T^{(Mql)}$-linear and
$\coker\Psi$ is a $\T^{(Mql)}$-module. \stpmc{} (Brown, as in Lemma~\ref{lem:s03-dcoset}).
\end{lemma}
\begin{proof}
$c\,T_\beta^\Gamma c^{-1}=T^{\Gamma'}_{\alpha'^{-1}\beta\alpha'}$ by the compatibility of $c_g$ with $\mathrm{res}$ and $\mathrm{cor}$ (proof of Lemma~\ref{lem:s03-dcoset}(a)), and
$\alpha'^{-1}\beta\alpha'=\beta$ (both diagonal). For the restriction, let $\Delta\in\{\Gamma,\Gamma'\}$ and $H:=\Gamma_0(Mq)\subseteq\Delta$. For $h=\left(\begin{smallmatrix}a&b\\c&d\end{smallmatrix}\right)\in H$,
$\beta^{-1}h\beta=\left(\begin{smallmatrix}a&b/p\\pc&d\end{smallmatrix}\right)$; it lies in $\Delta$ iff $p\mid b$ ($p\ne q$), and then it lies in $H$. So
$H\cap\beta\Delta\beta^{-1}=H\cap\beta H\beta^{-1}=\{h\in H:p\mid b\}$, and $H/(H\cap\beta H\beta^{-1})\to\Delta/(\Delta\cap\beta\Delta\beta^{-1})$ is injective. Both indices are $p+1$:
$H$ and $\Gamma$ map onto $\SL_2(\F_p)$ (as in Lemma~\ref{lem:s03-groups}(b), since $p\nmid Mq$), the groups $\{p\mid b\}$ are the preimages of the lower triangular group, and
$\Gamma'=\alpha'^{-1}\Gamma\alpha'$ with $\alpha'$ commuting with $\beta$. So the map is bijective, and we may take the representatives $r_i$ of
Lemma~\ref{lem:s03-Tclassical}(a) in $H$, for $\Delta$ and for $H$ simultaneously. For $h\in H$ the index $j(i)$ is then the same for both, and
$r_i^{-1}hr_{j(i)}\in H\cap\beta H\beta^{-1}$, so the cocycle formulas for $(T^\Delta_pz)|_H$ and $T^H_p(z|_H)$ agree term by term. Then
$T^{H}_p\Psi(x,y)=\mathrm{res}\,T^\Gamma_px-\mathrm{res}'\,T^{\Gamma'}_pcy=\mathrm{res}\,T^\Gamma_px-\mathrm{res}'\,cT^\Gamma_py=\Psi(T^\Gamma_px,T^\Gamma_py)$.
\end{proof}

\begin{lemma}[$H^1(G,V)=0$]\label{lem:s03-H1G}
For $k\ge4$ even, $l>k$, $l\ne q$: $H^1(G,V)=0$, i.e.\ $\Psi$ is injective. \stpmc{} (Ramanujan--Petersson [De74, Thm.~8.2] \stcited{}, not in
\pathtt{literature/}; $T_q$ on all of $H^1(\Gamma,V\otimes\C)$, Eisenstein part included, by Lemma~\ref{lem:s03-Tclassical}(d),(e); the Hecke eigenvalues of Eisenstein series
\lit{stein\_2007\_modular-forms-a-computational-approach\_draft}{Thms.~5.9, 5.10, p.~88} \stcited).
\end{lemma}
\begin{proof}
By Prop.~\ref{prop:s03-MV}, $H^1(G,V)\cong K:=\ker\Psi\subseteq H^1(\Gamma,V)^2$. Clearly $K\subseteq\ker(\Psi^*\Psi)$, and as the entries of $\Psi^*\Psi$ lie in the commutative
ring $\Z_{(l)}[T_q]$, $\mathrm{adj}(\Psi^*\Psi)\Psi^*\Psi=\Delta\cdot1$ with $\Delta=\det(\Psi^*\Psi)=-q^{2-k}P(T_q)$, $P(X):=X^2-q^{k-2}(q+1)^2$ (Prop.~\ref{prop:s03-PsiPsi}). So
$K\subseteq\ker(P(T_q))^2$. $H^1(\Gamma,V)$ is torsion-free and finitely generated over $\Z_{(l)}$ (Prop.~\ref{prop:s03-MV}), so it embeds into
$H^1(\Gamma,V)\otimes\C=H^1(\Gamma,V\otimes\C)$ (flatness, \lit{hida\_1993\_elementary-theory-of-l-functions-and-eisenstein-series}{\S6.2, (1a), p.~168}). It suffices
that no generalised eigenvalue $\lambda$ of $T_q$ on $H^1(\Gamma,V\otimes\C)$ satisfies $\lambda^2=q^{k-2}(q+1)^2$. By Lemma~\ref{lem:s03-Tclassical}(d),(e)
these eigenvalues are those of the classical $T_q$ on $S_k(\Gamma)$, on $\overline{S_k(\Gamma)}$ and on the Eisenstein space $E_k(\Gamma)$.
\emph{Cusp forms:} $\lambda=a_q(f)$ or $\overline{a_q(f)}$ for an eigenform $f$ of level dividing $M$, and $|\lambda|\le2q^{(k-1)/2}<(q+1)q^{(k-2)/2}$, as
$q+1-2\sqrt q=(\sqrt q-1)^2>0$. \emph{Eisenstein series:} $E_k(\Gamma)$ has the basis $E_{k,\chi,\psi}(q^t)$, $RLt\mid M$, $\chi\psi=1$ (Stein, Thm.~5.9); $T_q$ commutes with $f(w)\mapsto f(tw)$ for $q\nmid t$
(from $a_n(T_qf)=a_{qn}(f)+q^{k-1}a_{n/q}(f)$, \lit{stein\_2007\_modular-forms-a-computational-approach\_draft}{proof of Prop.~9.10, p.~164} \stcited), so
$\lambda=\chi(q)+\psi(q)q^{k-1}$ (Thm.~5.10) with $|\chi(q)|=|\psi(q)|=1$ ($q\nmid M$), so
$|\lambda|\ge q^{k-1}-1\ge q^{k/2+1}-1$ (as $k\ge4$), and $q^{k/2+1}-1-(q+1)q^{(k-2)/2}=q^{k/2-1}(q^2-q-1)-1\ge q-1>0$. Hence $P(T_q)$ is injective and $K=0$.
(Consistent with \pathtt{code/results/s03/betti\_det.json}: $\det(\Psi^*\Psi)\ne0$ on $S_k$ in all $121$ cases, and the HACS data of Lemma~\ref{lem:s03-Tclassical}, \stcomp.)
\end{proof}

\begin{definition}[Betti level-raising module]\label{def:s03-LR}
Let $\T$ be the (commutative, $\Z_{(l)}$-finite) Hecke algebra of $\Gamma_0(M)$ acting on $H^1(\Gamma,V)$; it contains $T_q=T_{\alpha'}$ and the $T_p^\Gamma$ of
Lemma~\ref{lem:s03-PsiHecke}. The entries of $\Psi^*\Psi$ lie in $\T$, so $\coker(\Psi^*\Psi)$ is a $\T$-module (unlike $\coker\Psi$, on which only
$\T^{(Mql)}$ acts). For a maximal ideal $\m$ of $\T$ put
$\mathrm{LR}_q(M,k)_\m:=\coker(\Psi^*\Psi)_\m$. It is a finite module whenever $T_q^2-q^{k-2}(q+1)^2$ is a non-zero-divisor on
$H^1(\Gamma,V)_\m$, which always holds under the hypotheses of Lemma~\ref{lem:s03-H1G} (its proof shows that $T_q^2-q^{k-2}(q+1)^2$ is injective on
$H^1(\Gamma,V)$). \status{definition} It is used by stage 06 as the Betti incarnation of $\cc_q$ (seed C3, T4; status there
\stsketch). The classical source for level raising is [Ri84]
\lit{ribet\_1984\_congruence-relations-between-modular-forms\_scan}{\S2, (2.2), p.~506; \S3, Prop.~3.3 and the matrix $a'\circ a$, p.~509}
\stcited{} (weight $2$; condition $a_q^2\equiv(1+q)^2$, matrix $\left(\begin{smallmatrix}q+1&\tau\\\tau&q+1\end{smallmatrix}\right)$; image-only scan, read on the page images).
\emph{Attribution fixed:} [Ri83] is \emph{not} a level-raising paper. It concerns congruence primes: [Ri83, Thm~(1.4), p.~195]
\lit{ribet\_1983\_mod-p-hecke-operators-and-congruences-between-modular-forms}{Thm (1.4), p.~195} \stcited{} says that a large prime $p$ (prime to the level, $p\ge k$) is a
congruence prime for $f$ iff $p$ divides the cohomological index $(\tilde L:L)$; this is the cohomological congruence module of a single form, an analogue of
$\cc_q$ but not the old/new module of level raising. The \pathtt{SPLIT.md} line ``\texttt{ribet\_1983} (level raising)'' and the stage-tex sentence ``the
cohomological old/new congruence module also appears in [Ri83]'' are corrected accordingly.
\end{definition}
\begin{proof}[Proof of the finiteness claim]
Let $\T$ be the ($\Z_{(l)}$-finite, commutative) Hecke algebra, $H:=H^1(\Gamma,V)_\m$, a finitely generated module over the noetherian local ring $\T_\m$ of Krull dimension $\le1$.
The entries of $\Psi^*\Psi$ lie in the commutative ring $\T_\m$ ($T_q\in\T$), so with the adjugate matrix $\Psi^*\Psi\cdot\mathrm{adj}(\Psi^*\Psi)=\Delta\cdot1$, $\Delta:=\det(\Psi^*\Psi)$, and
$\mathrm{im}(\Psi^*\Psi)\supseteq\Delta H^2$. So $\mathrm{LR}_q(M,k)_\m$ is a quotient of $(H/\Delta H)^2$. By Prop.~\ref{prop:s03-PsiPsi},
$\Delta=-q^{2-k}(T_q^2-q^{k-2}(q+1)^2)$ with $q$ a unit; if it is a non-zero-divisor on $H$, then $\dim H/\Delta H\le\dim H-1\le0$, so $H/\Delta H$ has finite length over
$\T_\m$, whose residue field is finite. Hence $\mathrm{LR}_q(M,k)_\m$ is finite. \stproved
\end{proof}

\begin{conjecture}[Ihara lemma for $\GL_2$, Betti form]\label{conj:s03-iharaGL2}
Let $\m$ be a maximal ideal of $\T^{(Mql)}$ (Lemma~\ref{lem:s03-PsiHecke}); call $\m$ \emph{Eisenstein} if there is $d\ge1$ such that $T_p-2\in\m$ for all but finitely
many primes $p\equiv1\bmod d$. This is the notion of [DT, p.~439--440]
\lit{diamond-taylor\_1994\_non-optimal-levels-of-mod-l-modular-representations}{p.~439--440}; DT's second condition $S_p-1\in\m$ is automatic for $p\equiv1\bmod dl$, since the
scalar $p$ acts as $p^{k-2}$ (Lemma~\ref{lem:s03-dcoset}(b)) and $l\in\m$. Compare Cor.~\ref{cor:s03-eis} on the quaternionic side.
For non-Eisenstein $\m$, $(\coker\Psi)_{\m,\mathrm{tors}}=0$. Equivalently, $\Psi\otimes\F_l$ is injective after localisation at
$\m$, because the groups $H^1(\cdot,V)$ are torsion-free (so the Tor sequence of $\otimes\F_l$ has no extra term; $H^2(\Gamma,V/l)$ plays no role) and $\Psi$ is injective (Lemma~\ref{lem:s03-H1G}). \stexp{} (consistent with C3: all torsion found is
Eisenstein, at $l=k+1$). The classical route uses Ihara's lemma [Ih75] (not in \pathtt{literature/}; \texttt{ihara\_1966} there is a different paper) or the congruence subgroup
property of $\SL_2(\Z[\frac1q])$ (\pathtt{literature/serre\_1970\_*}, scan, theorem not checked). The indefinite quaternionic analogue
is [MS19, Thm~1.1]
\lit{manning-shotton\_2019\_iharas-lemma-for-shimura-curves-over-totally-real-fields-via-patching\_arxiv}{Thm 1.1, p.~2} \stcited.
The ``equivalently'' needs $\ker\Psi=H^1(G,V)=0$, which is Lemma~\ref{lem:s03-H1G}. Writing $0\to K\to A\xrightarrow\Psi B\to C\to0$ with
$A=H^1(\Gamma,V)^2$, $B=H^1(\Gamma_0(Mq),V)$ free over $\Z_{(l)}$ (Prop.~\ref{prop:s03-MV}), one has $\ker(\Psi\otimes\F_l)\supseteq K/l$ and
$\ker(\Psi\otimes\F_l)/(K/l)\cong C[l]$; so, as $K=0$, $\Psi\otimes\F_l$ is injective after localisation at $\m$ iff $C[l]_\m=0$ iff $C_{\m,\mathrm{tors}}=0$
(all maps are $\T^{(Mql)}$-linear by Lemma~\ref{lem:s03-PsiHecke}).
\end{conjecture}

\begin{interface}[\pathtt{if:s03-amalgam}]\label{if:s03-amalgam}
\emph{Hypotheses:} $(M,q)=1$, $l$ prime with $l>k$ and $l\ne q$ (so $l\ge5$; $l\mid M$ is allowed), $k\ge4$ even, $V=\Sym^{k-2}(\Z_{(l)}^2)$ of the standard (column)
representation. $k=2$ is excluded: then $H^0\ne0$ (for $k=2$ the sequence stays exact on the left, since $H^0(\Gamma,V)^2\to H^0(\Gamma_0(Mq),V)$ is onto, but it does
not have the form below).
\emph{Exports:} Prop.~\ref{prop:s03-amalgam} (amalgam for $\Gamma_0(M)[\frac1q]$), Prop.~\ref{prop:s03-MV} (MV sequence, $\Psi$; with
Lemma~\ref{lem:s03-H1G}, $0\to H^1(\Gamma,V)^2\xrightarrow\Psi H^1(\Gamma_0(Mq),V)\to H^2(G,V)\to0$),
Prop.~\ref{prop:s03-PsiPsi} ($\Psi^*\Psi=\left(\begin{smallmatrix}q+1&-q^{2-k}T_q\\-T_q&q+1\end{smallmatrix}\right)$, with
$\det=(q+1)^2-q^{2-k}T_q^2$), Lemma~\ref{lem:s03-PsiHecke} ($\Psi$ is $\T^{(Mql)}$-linear), and Def.~\ref{def:s03-LR} ($\mathrm{LR}_q(M,k)_\m=\coker(\Psi^*\Psi)_\m$).
\emph{Status:} amalgam, MV, $\Psi^*\Psi$, Hecke equivariance and $H^1(G,V)=0$ \stpmc. $T_q=T_{\alpha'}$ is the classical operator: \stpmc{}, on $H^1_P$ by
[Hi93, \S6.3, Thm.~1] (Lemma~\ref{lem:s03-Tclassical}(d)) and on the Eisenstein part by Lemma~\ref{lem:s03-Tclassical}(e); \stcomp{} on all of $H^1$ for the $14$ HACS
cases of \pathtt{code/results/s03/C3\_rerun.json} (commit \texttt{05f7fe72}), and $\det(\Psi^*\Psi)\ne0$ on $S_k$ for the $121$ cases of
\pathtt{code/results/s03/betti\_det.json} (commit \texttt{0f3bd5e1}). Downstream should use only $\det$ and $\coker$ (D4).
\end{interface}
